\documentclass[11pt]{article}
\usepackage{preamble}
\renewcommand{\answer}[1]{}

\begin{document}

\ladderheading{AP Statistics \textperiodcentered\ Unit 2 \textperiodcentered\ Topic 2.11 \textperiodcentered\ B: 2026-11-10 \textperiodcentered\ E: 2026-12-04}{Comparing a Score of 92}

\focusbanner{TODAY'S PROBLEM}

Suppose exam scores at School A and School B are each approximately normal. Both distributions have mean $80$ points, but School A has standard deviation $9$ points and School B has standard deviation $15$ points.\par\smallskip A counselor says, ``A score of 92 is at about the 90th percentile at both schools because both schools have the same mean.'' A data coach says, ``The different spreads could put 92 at different percentiles.''\par
\begin{center}
\begin{tikzpicture}
\begin{axis}[
  width=0.88\linewidth,
  height=1.75in,
  xmin=35, xmax=125,
  ymin=0, ymax=0.050,
  axis x line=bottom,
  axis y line=left,
  ytick=\empty,
  xtick={50,65,80,92,110},
  xlabel={Exam score (points)},
  tick label style={font=\small},
  label style={font=\small},
  legend style={font=\small,draw=none,fill=none,at={(0.5,1.02)},anchor=south}
]
\addplot[dayblue,very thick,domain=35:125,samples=181]
  {1/(9*sqrt(2*pi))*exp(-0.5*((x-80)/9)^2)};
\addplot[warmred,very thick,dashed,domain=35:125,samples=181]
  {1/(15*sqrt(2*pi))*exp(-0.5*((x-80)/15)^2)};
\addplot[black,densely dotted] coordinates {(92,0) (92,0.049)};
\legend{{School A: $\mu=80,\ \sigma=9$},{School B: $\mu=80,\ \sigma=15$}}
\end{axis}
\end{tikzpicture}
\end{center}
\begin{firsttakebox}{FIRST TAKE --- before the video (5 min)}
Would a 92 stand out equally at both schools? Give one reason from the picture.
\workspace{0.9in}
\end{firsttakebox}

\begin{callout}{\IconBook\ STANDARDIZE BEFORE COMPARING (keep on the page)}{calloutgreen}
For a normal random variable, $z=\dfrac{x-\mu}{\sigma}$ tells how many standard deviations $x$ is from its mean.
The percentile of $x$ is the percent of values at or below $x$.
Use a normal table or calculator to find the area to the left of the $z$-score; multiply that area by 100.
A larger percentile means a higher standing relative to that school's scores.
\end{callout}

\newpage
\textbf{Finish after the video:}\par\smallskip

\dokbadge{1}~\textbf{(a)} Identify which school has the more variable score distribution and state how the corresponding normal curve shows that difference.\par
\workspace{0.7in}

\dokbadge{2}~\textbf{(b)} For a score of 92, calculate the $z$-score and percentile at each school. Show the standardization and report each percentile to the nearest tenth of a percent.\par
\workspace{1.4in}

\dokbadge{3}~\textbf{(c)} In 3--4 sentences, decide whose claim is better supported. Use both percentiles from (b) and explain how the different spreads change what a score of 92 means at each school.\par
\workspace{2.0in}

\begin{sentenceframebox}\raggedright\textbf{Frame:} At School A, $z=\blank[0.7in]$ and the percentile is \blank[0.8in]; at School B, $z=\blank[0.7in]$ and the percentile is \blank[0.8in].\end{sentenceframebox}

\begin{sentenceframebox}\raggedright\textbf{Frame:} The counselor's claim is \blank[1.0in] because the same 12-point gap is \blank[2.4in]; therefore a 92 has the higher relative standing at \blank[0.8in].\end{sentenceframebox}

\begin{summaryexitbox}{EXIT --- turn this in}
One thing the video changed about my first take: \hrulefill\par
\workspace{0.4in}
\end{summaryexitbox}

\end{document}
